\section{Introduction}

Graph services increasingly evaluate many analytical queries over the same graph.  A routing service answers paths from different origins, a social-network backend evaluates reachability or influence for many users, and an analyst submits a window of the same graph algorithm with different sources or parameters.  We refer to this workload as \emph{concurrent graph query processing} (CGQ).  Unlike a single graph analytic, a CGQ workload contains repeated access to one largely static topology while retaining independent query states and results.  Its systems objective is therefore to complete the query window efficiently without changing any constituent query.

GPUs are attractive for CGQ because they provide high memory bandwidth and parallel processing capacity.  Yet graph computation is dominated by data movement: adjacency lists are irregular, active vertices change from one iteration to the next, and query states are scattered across the graph.  Concurrency does not remove this bottleneck.  When two queries reach the same vertex, independent execution accesses the same neighbor list twice; when their traversals diverge, executing them together can add contention without reuse.  In our pilot profiles, exposed data-access stalls account for 43.1--59.8\% of issue stalls in representative iterations even though achieved bandwidth ranges from only 31.0--62.0\% of peak.  These measurements point to memory-access latency, rather than insufficient arithmetic, as the central bottleneck for our target workload.

Recent database graph systems address important but different limitations.  TGraph provides a tensor-centric model that hides accelerator-specific details~\cite{zhang_tgraph_2025}; CGgraph improves out-of-memory processing and CPU--GPU cooperation~\cite{cui_cggraph_2024}; and CompressGraph reuses repeated neighbor sequences through direct processing on a compressed representation~\cite{chen_compressgraph_2023}.  These systems primarily optimize one graph analytic at a time.  Concurrent graph systems instead share computation or coordinate multiple queries: \textsc{iBFS} fuses BFS instances on GPUs~\cite{liu_ibfs_2016}, while \textsc{Glign} groups and delays traversals to align expensive phases~\cite{yin_glign_2022}.  Apart from specialized BFS execution, however, existing work does not systematically organize general CGQ execution around its memory-access bottleneck.

Our key observation is that concurrency creates both the problem and the opportunity.  It increases the number of active states, but it also exposes repeated neighborhood accesses and a new regular dimension across queries.  We can therefore mitigate memory-access overhead in two complementary ways: \emph{reduce the number of accesses} by sharing a neighborhood among exactly the queries that currently need it, and \emph{reduce the effective cost of the remaining accesses} by processing query states in a regular order.  Query scheduling provides a third lever: it changes when traversals reach the same graph region and hence how much access reuse is available.

Figure~\ref{fig:core-idea} summarizes this principle: \emph{share repeated accesses and regularize the remainder}.  This is a workload-level view rather than a GPU execution detail.  It first asks which graph data and query states are logically reusable, then maps the resulting computation to the hardware in Section~\ref{sec:implementation}.

\begin{figure}[t]
  \centering
  \begin{tikzpicture}[
      font=\scriptsize,
      box/.style={draw, rounded corners, align=center, minimum height=6mm},
      arrow/.style={-{Latex[length=1.5mm]}, thick},
      cost/.style={font=\scriptsize, align=center}
    ]
    \node[box, fill=gray!8] (i0) at (0,0.7) {$q_0$ frontier};
    \node[box, fill=gray!8] (i1) at (0,-0.05) {$q_1$ frontier};
    \node[box, fill=red!8, minimum width=16mm] (r0) at (1.65,0.7) {CSR row $v$};
    \node[box, fill=red!8, minimum width=16mm] (r1) at (1.65,-0.05) {CSR row $v$};
    \draw[arrow] (i0) -- (r0);
    \draw[arrow] (i1) -- (r1);
    \node[cost] at (0.8,-0.65) {independent: two neighborhood accesses};

    \draw[dashed, gray] (2.65,-0.75) -- (2.65,1.05);

    \node[box, fill=blue!8] (mask) at (3.65,0.7) {$v: \querymask(v)=0011_2$};
    \node[box, fill=green!8] (row) at (5.25,0.7) {one shared access};
    \node[box, fill=blue!8] (state) at (4.45,-0.05) {$x(v,0{:}Q{-}1)$};
    \draw[arrow] (mask) -- (row);
    \draw[arrow] (state) -- (row);
    \node[cost] at (4.45,-0.65) {\system: share repeats; regularize state};
  \end{tikzpicture}
  \caption{The core principle: share repeated accesses and regularize the remainder.  \system represents exact query overlap at a vertex, accesses its neighborhood once, and keeps the associated query states together.}
  \Description{Two independent query frontiers each access the same adjacency row. GraphWeft combines them into one vertex with an exact query mask, one shared neighborhood access, and a contiguous row of query values.}
  \label{fig:core-idea}
\end{figure}


Realizing this opportunity raises three challenges.

\paragraph{C1: Reuse must preserve query semantics.}
Two queries may visit the same vertex at different logical levels or with different values.  A query-oblivious union can access a neighbor list once but may apply unrelated queries to it, replacing redundant memory accesses with false computation.  The system must expose exact, iteration-local overlap while preserving independent values and convergence state.

\paragraph{C2: Access reduction and access regularity have different win regions.}
Sharing an active vertex reduces neighborhood accesses when frontiers are sparse and overlapping.  During broad phases, however, irregular updates to many query states can dominate, and a destination-oriented traversal with regular query-state access may be preferable even if it examines more edges.  No single traversal strategy minimizes memory overhead across all graphs and iterations.

\paragraph{C3: Shareability changes over time.}
Queries that eventually visit similar regions may reach them in different iterations.  Grouping and start-time alignment can increase simultaneous overlap, but excessive delay or planning cost can erase the benefit.  Coordination must therefore improve realized access reuse rather than optimize a static source-similarity score alone.

We present \system, a CGQ engine for a GPU-resident graph.  Its thesis is that \emph{CGQ performance can be improved by treating memory-access overhead as a unified design objective across query representation, graph traversal, and scheduling}.  \system uses exact per-vertex query masks to represent current overlap.  \sharedexpand shares each active vertex's neighborhood among only the queries named by its mask, reducing topology accesses.  \coalescedgather processes the remaining query states in a vertex-major order, reducing their effective access cost during broad phases.  A phase-aware scheduler groups and offsets queries to expose more simultaneous vertex overlap, while an adaptive selector chooses the traversal strategy from exact runtime state.  Approximate scheduling may miss an optimization opportunity, but it cannot change a query answer.

This paper makes four contributions:

\begin{itemize}
  \item We formulate GPU-based CGQ as a memory-access problem and decompose its overhead into repeated topology accesses and the effective cost of query-state accesses.  Pilot measurements connect both terms to observed execution behavior (Section~\ref{sec:background}).
  \item We develop an exact vertex-sharing computation model and two complementary traversal strategies: \sharedexpand reduces repeated neighborhood accesses, while \coalescedgather organizes broad query-state aggregation for efficient access.  Both support typed BFS, SSSP, and widest-path state rather than only BFS status bits (Section~\ref{sec:design}).
  \item We design a sharing-aware scheduler that estimates pairwise traversal phases, forms query batches, and selects bounded start offsets to increase realized vertex sharing.  Runtime decisions use exact query state, so prediction quality affects performance but not correctness (Sections~\ref{sec:overview} and~\ref{sec:planner}).
  \item We implement these ideas in a 7.6-KLoC C++17/CUDA prototype and define an evaluation that measures end-to-end CGQ throughput, memory-access reduction, component contributions, scheduling quality, scalability, and overhead (Sections~\ref{sec:implementation} and~\ref{sec:evaluation}).
\end{itemize}

Current results are preliminary and delimit the remaining work.  For all-push BFS with 128 queries and 32 resident slots, online grouping and offsets reduce union-edge accesses by 15.5--64.4\% and improve batch execution by $1.15\times$--$2.31\times$ on six graphs.  The current fixed-$Q{=}64$ hybrid path, however, reaches only $0.49\times$ the geometric-mean throughput of \textsc{iBFS} on four graphs.  Section~\ref{sec:evaluation} therefore treats fixed-policy performance, the traversal selector, and the same-engine \textsc{iBFS}/\textsc{Glign} decomposition as gating evidence rather than hiding the gap behind favorable configurations.
